
In order for robots to perform daily tasks in their environment, it is essential to have an accurate representation of a scene and relevant objects. One prominent explicit 3D representation is Gaussian Splatting (3DGS) \cite{3dgs}, which reconstructs a 3D scene from RGB images and camera poses. 3DGS has been applied to a plethora of robotic tasks including manipulation \cite{lu2024manigaussian,shorinwa2024splat}, SLAM \cite{monoGS, yan2023gs} and navigation \cite{chen2024splatnav,lei2024gaussnav}. The fast training speeds and high accuracy reconstruction of 3DGS it an effective representation for robots; however, when provided with a limited number of views of a scene or object quality quickly deteriorates \cite{swann2024touch}. 

Traditionally, the creation of 3DGS scenes is done with a). many views and b). complete human supervision, where visual inspection allows for a person to manually choose views to take before training a 3DGS. 3DGS is well-known for requiring tens of views and near-perfect camera poses; as it is prone to overfitting on sparse input views.  This approach, however, is inconsistent with many of the potential robotic applications of 3DGS. In a robotic setting every additional view is costly, requiring movement and potentially limited onboard compute. It is often inefficient to collect hundreds of views required for training a 3DGS model in robotic environments. For future applications of 3DGS on robotics systems, robust few-view training methods are needed.
%Forms of depth regularization have emerged as methods to unify Gaussians from different views but the same objects \cite{chung2024depthregularized}. However, these methods often make assumptions and operate of high quality visual data. For computational efficiency, a robot may have a low resolution camera and depth image, making few view GS more challenging. 

We address this issue from two fronts. First we propose \textit{semantic depth alignment} to enhance the impact of each additional view. We draw from the geometric power of state of the art monocular depth estimation models, and with an off-the-shelf depth camera and color image, we perform a semantic depth alignment by aligning objects and the background of an image with Segment Anything Model 2 (SAM2) \cite{ravi2024sam2segmentimages}. This is used as a form of initialization in our few shot scene. During training, depth supervision is enforced with a relaxed relative loss, which mitigates the scale ambiguities between the true depth and estimated depth from the depth estimation models. Gentle normal supervision and camera optimization guides the splat to reconstruct a visually accurate model.
% We first address this by improving few shot Gaussian Splats trained in the wild.
\begin{figure}[t]
    \centering
    \includegraphics[width=.45\textwidth]{figure/splash_next_best.png}
    \caption{Our method outperforms random view selection for few view 3DGS scenes. We show a series of robot views, with the next view proposed by our method compared to random. Our method maximizes the Fisher information gain for each view over both vision and touch inputs.}
    \label{fig:showcase}
    % We also highlight our method versus 3DGS with monocular depth supervision.
    \vspace{-6mm}
\end{figure}

% part 3: autonomously with a robot -- which views
In addition to maximizing the impact of each of the limited training views, we propose a method to optimally select the \emph{next best view} based on depth uncertainty, and optionally color uncertainty. Even when there are a limited number of views available, uncertainty naturally emerges from depth discontinuities, occluded regions, and sources of aleatoric error. It is critical to understand where \textit{next} best views and even touches should be made in a scene to reduce this uncertainty and construct precise scenes that robots can effectively operate in. We draw from recent works in next best view selection for Gaussian Splatting, which use the expected Fisher Information gain of candidate views \cite{jiang2023fisherrf}, and extend it to select views \textit{and} touches based on \textbf{depth uncertainty}. We posit that the geometry and visual quality of a scene can be guided through depth, as regions of erroneous color are often associated with ambiguous depth. This framework can be easily extended to touch, which has been shown to improve 3DGS quality \cite{swann2024touch}, as we can use the local depth uncertainty of candidate touches to propose where to touch next. We finally present a method to supervise on singular touches by adding local cameras to the 3DGS training set that only trains on depth. Altogether, we construct a framework that enables for the autonomous sense selection with a robotic manipulator, which we call \textit{Next Best Sense}.

% . \gls*{NeRF}s have been applied to a number of robotics challenges, including path planning \cite{chen2023catnips} and manipulation \cite{suresh2023neural}. The method of \gls*{3DGS} \cite{3dgs} has recently advanced the field by providing high-quality and high-speed training along with real-time rendering. Precise representations and real-time visual reconstruction are important aspects for robots to interact with their environments. However, in most cases, visual-only information is not sufficient to interact with complex objects. Situations as common as reaching into a cluttered drawer or grasping something in the dark necessitate a sense of touch. 




\subsection{Key Contributions}
We propose Next Best Sense, the first approach that enables uncertainty-guided view \textit{and} touch selection for training a \gls*{3DGS} for robotic manipulators. Our main contributions in Next Best Sense are as follows:
\begin{itemize}
    \item We propose an improved method in few-shot 3DGS scenes with a novel depth alignment method using Segment Anything Model 2.
    \item We present a novel extension of FisherRF by leveraging depth uncertainty to assist in view and touch selection.
    \item We present a closed-loop framework coupling both perception and action for training few-view scenes for robotic manipulators with FisherRF. To the best of our knowledge, this is the first work that enables next best view planning for Gaussian Splatting for robot manipulators in real-time. 
    \item We extend our method to uncertainty-guided touch, demonstrating qualitative improvements with only 10 touches.
\end{itemize}

The remainder of the paper is outlined as follows. Section \ref{sec:related-works} first covers existing literature in the field. Section \ref{sec:prelim} covers mathematical preliminaries, and Section \ref{sec:method} introduces our method. Section \ref{sec:results} showcases our results, and finally, Section \ref{sec:conclusions} discusses these results and future avenues of research.


\section{Related Works}
\label{sec:related-works}
\textbf{Few Shot Gaussian Splatting}. Traditional Gaussian Splats trained on few, sparse views tend to overfit to each view, causing degradation on novel views. To address this, prior works  enforce monocular depth priors \cite{chung2024depthregularized, paliwal2024coherentgssparsenovelview}, use vision priors such as DUSt3R \cite{dust3r2023}, and unpooling of sparse Gaussians from SfM points \cite{zhu2023fsgs} to generalize to novel views. These works often rely on simple alignment or priors for accurate depth information, which is challenging in robotics, where alignment is performed with noisy depth data, and the views are even more sparse.

\textbf{Fusing Vision and Touch for Robotics.} Prior works have addressed the multi-modal fusion of vision and touch to reconstruct accurate objects for manipulation \cite{smith2021active, suresh2023neural}. Others have focused on constructing a shared representation for manipulation \cite{yang2024binding, fu2024touch}.  \cite{watkins2019multi, smith20203d} focus on many views or only tactile input. Recently, there are also works that try to incorporate tactile data into Gaussian Splatting radiance field. Touch-GS \cite{swann2024touch} and Snap-It, Tap-It, Splat-It \cite{comi2024snap} address tactile data for \gls*{3DGS}. While Touch-GS demonstrates the improvements on GS scenes with touch, it suffers from two shortcomings: 1. it requires a human to move the robot to desired vision \textit{and} touch poses and 2. requires the order of 100-500 touches to construct a visually and geometrically sufficient scene. \cite{comi2024snap} focuses on all manually selected several views and touches in the real world.

% \textbf{Uncertainty in Radiance Fields.} Uncertainty quantification in radiance fields is a relatively new area of study. Uncertainty has traditionally been applied to implicit representation such as NeRF, which require learning complex additions onto the traditional pipeline \cite{pan2022activenerf, shen2021snerf} or simply using the entropy of the ray distribution to measure uncertainty. Bayes' Rays \cite{goli2023} utilizes the Fisher information on a spatial deformation field on NeRF by measuring color uncertainty. Uncertainty in Gaussian Splatting remains generally unexplored with using variational inference to quantify uncertainty. In general, these methods use \textit{color} as the guiding metric for uncertainty.

% I think we can combine these two sections if we are low on space

\textbf{Active Learning for Radiance Fields.} The problem of active learning or next best view has been extensively studied by the vision and robotics community.
% Traditionally, frontier-based approach~\cite{frontier} and Rapidly exploring Random Tree~(RRT)~\cite{RRT} are the most representative and have been extended and enhanced for different data representations~\cite{zhu20153d, karaman2010incremental, dornhege2013frontier, shen2012autonomous} and complex conditions~\cite{ye2022multi, papachristos2017uncertainty}. 
Recently, as the radiance field gained popularity, there are also methods that try to find the next best view for the current model \cite{lin2022active}. Jin  trains a neural network to predict per-pixel uncertainty for view selection. However, their method requires an image capture as the input of the network, which is closer to dataset sub-sampling rather than active learning.
ActiveNeRF~\cite{pan2022activenerf} studied the next best view selection by directly adding variances as the output of the vanilla NeRF, and it is the closest approach to the view selection problem we present in this paper. Most recently, FisherRF~\cite{jiang2023fisherrf} quantized the expected Fisher Information Gain of candidate views on the radiance field, and selects views based on it. It was built on top of Gaussian Splatting~\cite{3dgs}.

\textbf{Next Best View for Robotic Applications in Radiance Fields.} Recent research on NBV selection traditionally uses the uncertainty of a neural radiance field via ray entropy \cite{lee2022uncertainty} and interpreting the radiance field as a probability distribution \cite{shen2024estimating}. Similar to our work, a first work of performing active planning in Gaussian Splatting \cite{jin2024gs} uses information gain; however, they are focused on the task of navigation on completely unseen areas.
